% !TeX program = xelatex
% !TEX TS-program = xelatex
% Compile with XeLaTeX or LuaLaTeX, twice for cross-references.
% Times New Roman where available; Times-compatible TeX Gyre Termes otherwise.
% No platform-specific font paths or redistributed Windows fonts are required.
% Existing images are referenced without editing or copying source assets.
\documentclass[11pt,a4paper]{article}
\usepackage[margin=22mm]{geometry}
\usepackage{fontspec}
% Retain Times New Roman where available; use Times-compatible Termes otherwise.
\IfFontExistsTF{Times New Roman}{
  \setmainfont{Times New Roman}
  \setsansfont{Times New Roman}
  \setmonofont{Times New Roman}
}{
  \setmainfont{TeX Gyre Termes}
  \setsansfont{TeX Gyre Termes}
  \setmonofont{TeX Gyre Termes}
}
\usepackage{amsmath,amssymb}
\usepackage{graphicx,float,caption,subcaption}
\usepackage{longtable,array,booktabs}
\usepackage{hyperref}
\hypersetup{hidelinks}
\usepackage{microtype}
\setlength{\emergencystretch}{3em}
\setlength{\parindent}{0pt}
\setlength{\parskip}{6pt}
\captionsetup{font=small,labelfont=bf}
% Paths are relative to the repository/Overleaf project root.
% Compile from that root; a missing image is an error, not a placeholder.
\newcommand{\projectimage}[2]{%
  \includegraphics[width=\linewidth,height=#2,keepaspectratio]{#1}%
}
\begin{document}
{\centering\Large\bfseries Stage-Resolved Transcriptomic Analysis Reveals Distinct Molecular Programs Associated with Liver Fibrosis Progression\par}



\section*{Abstract}
Understanding the molecular changes associated with progressive liver fibrosis is essential for characterizing disease heterogeneity and identifying biologically relevant markers. However, the relationship between fibrosis stage, transcriptional alterations, and their underlying biological processes remains incompletely characterized across independent cohorts. This study aimed to identify reproducible, stage-associated transcriptional patterns and characterize their biological organization through integrated transcriptomic and knowledge-graph analyses. We integrated eight publicly available human liver transcriptomic datasets comprising 1,085 samples from 1,027 recorded patients. Cross-study association analysis identified 3,845 fibrosis-associated genes and 4,199 genes associated with ordered disease-group labels, including 1,822 shared genes. Clustering of fibrosis-associated expression profiles across stages F0–F4 revealed five distinct patterns, encompassing early and late increases, an early rise followed by a plateau, and intermediate- and late-stage decreases. Functional enrichment implicated extracellular matrix organization, collagen remodeling, integrin and PDGF signaling, immune-related processes, and other stage-associated biological programs. A multi-evidence prioritization framework integrating progression strength, expression patterns, stage-transition evidence, fibrosis association specificity, and functional annotations identified 447 Tier-1 genes and a representative shortlist of 23 candidates. Leave-one-study-out analysis retained statistical significance for 442 Tier-1 genes across all five refits. Knowledge-graph analysis further highlighted interconnected extracellular matrix, cell-adhesion, migration, and signaling annotations among prioritized genes. Collectively, these findings characterize heterogeneous transcriptional programs associated with liver fibrosis severity and provide a reproducible computational framework for prioritizing fibrosis-associated genes and investigating their biological context.
\section*{Introduction}

Liver fibrosis is a central feature of chronic liver injury and an important indicator of clinical outcomes in fatty liver disease \cite{ref1}. In a prospective study of adults with non-alcoholic fatty liver disease (NAFLD), advanced fibrosis stages were associated with increased risks of liver-related complications and death \cite{ref2}. The disease spectrum has since been reclassified using the terms metabolic dysfunction-associated steatotic liver disease (MASLD) and metabolic dysfunction-associated steatohepatitis (MASH) \cite{ref3}. Nevertheless, fibrosis progression remains heterogeneous: a meta-analysis of paired-biopsy studies found that fibrosis may progress, remain stable, or regress over time \cite{singh2015}. Characterising the molecular changes associated with fibrosis severity is therefore important for understanding this heterogeneity and identifying candidate genes for further investigation.

Transcriptomic studies have identified gene-expression signatures associated with steatohepatitis and fibrosis, including signatures evaluated across independent replication cohorts \cite{govaere2020}. Multi-cohort analyses have further characterised expression signatures associated with MASLD and MASH \cite{piras2024}, shared and disease-specific fibrogenic programmes across liver disease aetiologies \cite{yang2026}, and transcriptional differences between advanced fibrosis (F3--F4) and mild or absent fibrosis (F0--F1) \cite{maiers2026}. Despite these advances, further work is needed to integrate cross-study evidence across the full fibrosis-stage spectrum (F0--F4), characterise distinct stage-associated expression patterns, and systematically connect these patterns to curated biological annotations.

Integrating public transcriptomic datasets presents several methodological challenges. Datasets differ in expression units, gene identifiers, clinical annotations, and representation of fibrosis stages. Many also retain historical NAFLD and non-alcoholic steatohepatitis (NASH) labels that cannot always be directly mapped to current diagnostic criteria. Repeated biopsies may lead to disproportionate patient representation if correlated observations are treated as independent, while unequal stage sizes can affect statistical comparisons. In addition, gene rankings and expression patterns may vary depending on the studies included in the analysis. Robust cross-study analysis therefore requires explicit data harmonisation, appropriate patient-level handling, and systematic assessment of the consistency and stability of the resulting associations.

Biomedical knowledge graphs offer a complementary approach to interpreting transcriptomic findings by representing relationships among genes, diseases, pathways, and other biological entities. Such resources have supported biomedical discovery \cite{himmelstein2017}. Emerging resources such as OptimusKG aim to integrate molecular, anatomical, clinical, and environmental information within a unified knowledge graph \cite{ref5}. Linking stage-associated expression patterns to curated graph evidence can provide biological context and support systematic candidate prioritisation. However, pathway enrichment and graph connectivity provide contextual evidence rather than proof of causal mechanisms or therapeutic efficacy.

The objective of this study was to identify reproducible, stage-associated transcriptional patterns in liver fibrosis and characterise their biological organisation through integrated transcriptomic and knowledge-graph analyses. We integrated eight public human liver transcriptomic datasets with curated OptimusKG evidence to address three research questions: (i) which gene associations with ordered fibrosis stages and disease-group labels are directionally consistent across studies; (ii) which expression patterns characterise the fibrosis-stage spectrum from F0 to F4; and (iii) how association strength, expression patterns, stage-transition evidence, and functional annotations can be combined to prioritise candidate genes and describe their biological relationships. To address these questions, we harmonised transcriptomic and clinical annotations, evaluated cross-study associations, characterised fibrosis-stage expression patterns, assessed internal robustness through leave-one-study-out analyses, and integrated prioritised genes with curated biological annotations.

The principal contributions of this study are as follows:

\begin{itemize}
    \item An integrated transcriptomic and knowledge-graph resource. Eight public datasets comprising 1,085 samples from 1,027 recorded patient IDs were harmonised and linked to OptimusKG. The fibrosis analysis included 610 independent staged patients, while the disease-group analysis included 216 patients. The integrated graph contained 126,244 nodes and 35,132,831 relationships.

    \item Cross-study associations and fibrosis-stage expression patterns. The analyses identified 3,845 fibrosis-associated genes and 4,199 genes associated with ordered disease-group labels, including 1,822 shared genes. Fibrosis-associated expression profiles were organised into five patterns across F0--F4, extending the analysis beyond comparisons restricted to broad disease categories or dichotomised fibrosis groups.

    \item Multi-evidence gene prioritisation. A framework combining association strength, expression patterns, stage-transition evidence, operational fibrosis-only status, and functional annotations identified 447 Tier 1 genes and a cluster-representative shortlist of 23 candidates, with up to five candidates selected per expression pattern.

    \item Internal robustness and biological-context assessment. 442 of the 447 Tier 1 genes retained statistical significance across all five leave-one-study-out refits. An enrichment-ranked subgraph connected 42 prioritised genes to twelve selected pathway or biological-process annotations through 154 relationships. These analyses assess internal robustness and the organisation of curated biological evidence rather than independent clinical or functional validation.
\end{itemize}

The remainder of this manuscript is organised as follows. The Results section presents the integrated resource, cross-study associations, fibrosis-stage expression patterns, functional enrichment, candidate prioritisation, robustness assessments and gene--annotation network. The Discussion interprets the findings in the context of prior work and considers their biological and methodological limitations. The Conclusion identifies the principal takeaways and future priorities. The Materials and Methods section then describes the data sources, harmonisation procedures, patient-handling rules, statistical analyses, knowledge-graph integration and gene-prioritisation framework.
\section*{Results}
The study workflow integrates transcriptomic harmonisation, cross-study association analysis and biological interpretation (Fig.~\ref{fig:6}).

\begin{figure}[H]\centering
\projectimage{manuscript_v3_pipeline.pdf}{0.80\textheight}
\caption{\textbf{Liver knowledge graph workflow.} Eight datasets undergo transcriptome processing, gene-identifier harmonisation, sample matching, clinical annotation and expression scaling before integration with OptimusKG. Association discovery informs disease-linked gene grouping, fibrosis-stage expression-pattern clustering, adjacent-stage expression comparisons, functional enrichment and multi-evidence candidate-gene prioritisation. Biological-network interpretation contextualises the prioritised genes.}\label{fig:6}
\end{figure}

\subsection*{Integrated resource and cross-study gene associations}
The integrated resource comprised eight human liver transcriptomic datasets containing 1,085 samples and 1,027 recorded patient identifiers (Table~\ref{tab:1}; workflow in Fig.~\ref{fig:6}). Five studies provided fibrosis-stage information for 668 biopsies, yielding 610 patient-level observations after repeat-biopsy handling. The ordered disease-group analysis included 216 patient-level observations across three studies. These populations were subsets of the full resource and differed between clinical axes. Patient identifiers were resolved within studies, while participant independence across accessions remained unestablished.

Gene harmonisation retained 53,993 measured genes, including a common set of 14,794 genes measured in all eight studies. Integration with OptimusKG produced a graph containing 126,244 nodes across eight entity types and 35,132,831 relationships across 14 types. Expression measurements contributed 32,471,042 relationships. The graph also contained 8,044 selected gene–disease progression links, comprising 3,845 fibrosis associations and 4,199 disease-group associations. Verification checks confirmed that these progression links matched the final analysis exports and that raw measurements and curated relationship counts were preserved. The knowledge layer retained broader biomedical coverage beyond liver-specific conditions. This resource provided the expression and annotation coverage for the subsequent association and prioritisation analyses.

Cross-study testing identified 3,845 fibrosis-associated genes and 4,199 genes associated with ordered disease-group labels, using the selection criteria described in Methods. The corresponding test families contained 14,794 and 14,784 genes. The selected sets overlapped at 1,822 genes and formed a union of 6,222 genes, with 2,023 fibrosis-only and 2,377 disease-group-only genes (Fig.~\ref{fig:2}a).

\clearpage
\begin{figure}[H]\centering
\begin{subfigure}[t]{0.98\textwidth}\centering
\projectimage{task1_venn.png}{0.32\textheight}\caption{}\end{subfigure}
\par\vspace{3pt}
\begin{subfigure}[t]{0.98\textwidth}\centering
\projectimage{task1_shared_rho_scatter.pdf}{0.40\textheight}\caption{}\end{subfigure}
\caption{\textbf{Overlap and direction of fibrosis and disease-group associations.} (a) The selected fibrosis and disease-group sets share 1,822 genes, with 2,023 fibrosis-only and 2,377 disease-group-only genes. (b) Median within-study correlations among shared genes; the displayed Spearman statistic \(\rho=0.702\), \(p=1.9\times10^{-270}\) compares the two correlation vectors across genes. These include 1,408 concordant positive, 384 concordant negative and 30 discordant associations. Selection required \(q<0.05\), agreement in direction across all testable studies and at least three contributing studies. Disease-group ordering represents source control/NAFL/NASH labels rather than a direct inflammation measurement.}\label{fig:2}
\end{figure}
\clearpage

The fibrosis-only set contained 1,040 positively and 983 negatively associated genes. Disease-group-only genes were predominantly positively associated, with 2,042 positive and 335 negative associations. Within the shared set, 1,408 genes had positive associations on both axes and 384 had negative associations on both; the remaining 30 were negatively associated with fibrosis but positively associated with disease-group ordering (Fig.~\ref{fig:2}b). The shared set therefore included both concordant and discordant associations. The axis-specific labels describe selection outcomes across different study sets rather than biological specificity. To examine how fibrosis-associated genes varied across severity, their F0--F4 expression profiles were analysed next.

\subsection*{Fibrosis-stage expression patterns and cluster stability}
The 3,845 fibrosis-associated genes formed five expression-pattern clusters (Fig.~\ref{fig:3}; Table~\ref{tab:2}). C1 contained 857 genes with an early increase followed by continued elevation, whereas C2 contained 1,367 genes with a late increase, largest between F3 and F4. C3 comprised 226 genes with an early rise followed by a plateau. The decreasing patterns comprised C4, with 397 genes showing a mid-stage decrease, and C5, with 998 genes showing a late decrease. C2 was the largest cluster and C3 the smallest. These profiles describe differences between patients at different stages rather than serial molecular trajectories within individuals.

\begingroup\fontsize{9}{11}\selectfont\setlength{\tabcolsep}{3pt}
\begin{longtable}{>{\raggedright\arraybackslash}p{0.065\textwidth}>{\raggedright\arraybackslash}p{0.31\textwidth}>{\raggedright\arraybackslash}p{0.06\textwidth}>{\raggedright\arraybackslash}p{0.1\textwidth}>{\raggedright\arraybackslash}p{0.055\textwidth}>{\raggedright\arraybackslash}p{0.07\textwidth}>{\raggedright\arraybackslash}p{0.07\textwidth}>{\raggedright\arraybackslash}p{0.07\textwidth}}
\caption{Cross-sectional clusters, confidence cores and prioritisation}\label{tab:2}\\
\toprule
\textbf{Cluster} & \textbf{Pattern} & \textbf{Genes} & \textbf{Confidence core\(\geq 0.8\)} & \textbf{Tier1} & \textbf{Shortlist} & \textbf{Mean F0} & \textbf{Mean F4} \\
\midrule\endfirsthead
\toprule
\textbf{Cluster} & \textbf{Pattern} & \textbf{Genes} & \textbf{Confidence core\(\geq 0.8\)} & \textbf{Tier1} & \textbf{Shortlist} & \textbf{Mean F0} & \textbf{Mean F4} \\
\midrule\endhead
\bottomrule\endfoot
C1 & Early increase (F0\ensuremath{\rightarrow}F1), then continued & 857 & 90 & 93 & 5 & -0.295 & 0.515 \\
C2 & Late increase (largest at F3\ensuremath{\rightarrow}F4) & 1367 & 806 & 262 & 5 & -0.227 & 0.86 \\
C3 & Early rise, then plateau & 226 & 74 & 7 & 5 & -0.27 & 0.136 \\
C4 & Mid decrease (largest at F1\ensuremath{\rightarrow}F2) & 397 & 126 & 3 & 3 & 0.161 & -0.279 \\
C5 & Late decrease (largest at F3\ensuremath{\rightarrow}F4) & 998 & 742 & 82 & 5 & 0.184 & -0.658 \\
\end{longtable}\endgroup

Stage means in Table~\ref{tab:2} are expressed on the original within-study standard-deviation scale; the clustering and profile panels use values additionally standardised across stages. The table also records the confidence-core and prioritisation counts to show how the different patterns contributed to subsequent analyses.

\clearpage
\begin{figure}[H]\centering
\begin{subfigure}[t]{0.98\linewidth}\centering
\projectimage{task2_cluster_trajectories.pdf}{0.73\textheight}\caption{}\end{subfigure}
\caption{\textbf{Five cross-sectional fibrosis expression patterns.} (a) Grey lines show individual gene profiles standardised across F0--F4; red lines show cluster centroids. (b) Heatmap of the same row-standardised profiles, grouped by cluster. (c) Cluster-number diagnostics showing silhouette scores and stability under gene subsampling and patient bootstrapping; error bars show standard deviations across 50 refits. Cluster sizes are 857, 1,367, 226, 397 and 998. These profiles represent cross-sectional stage associations rather than longitudinal patient trajectories; grey lines are individual genes, not confidence intervals.}\label{fig:3}
\end{figure}
\clearpage
\begin{figure}[H]\ContinuedFloat\centering
\setcounter{subfigure}{1}
\begin{subfigure}[t]{0.98\textwidth}\centering
\projectimage{task2_heatmap.pdf}{0.78\textheight}\caption{}\end{subfigure}
\caption{\textbf{Five cross-sectional fibrosis expression patterns (continued).} (b) Heatmap of the row-standardised gene profiles, with cluster sizes and the stage-profile colour scale.}
\end{figure}
\clearpage
\begin{figure}[H]\ContinuedFloat\centering
\setcounter{subfigure}{2}
\begin{subfigure}[t]{0.98\textwidth}\centering
\projectimage{task2_choosing_k.pdf}{0.76\textheight}\caption{}\end{subfigure}
\caption{\textbf{Five cross-sectional fibrosis expression patterns (continued).} (c) Cluster-number diagnostics. Error bars show standard deviations across 50 refits; the highlighted choice is k = 5.}
\end{figure}
\clearpage

The five-cluster solution met the implemented selection rule: it was the largest evaluated k with mean gene-subset ARI of at least 0.95. At \(k=5\), mean gene-subset ARI was 0.983 (SD 0.011), compared with patient-bootstrap ARI of 0.592 (SD 0.062). The silhouette score was 0.352; the highest silhouette score among evaluated solutions occurred at \(k=3\) (0.587). Thus, \(k=5\) reflected the stated stability rule rather than optimisation of the silhouette score. Study-adjusted profiles yielded cluster agreement of \(\mathrm{ARI}=0.901\) with the main analysis and a median per-gene profile correlation of 0.999.

In the additional centroid-matching bootstraps, median shape correlations ranged from 0.967 to 0.996, although minimum correlations were negative for C2, C3 and C5. A total of 1,838 genes met the conditional membership-confidence threshold of 0.8: 90 in C1, 806 in C2, 74 in C3, 126 in C4 and 742 in C5. Stable average shapes therefore coexisted with variable individual-gene membership. Adjacent-stage comparisons and enrichment analyses were used to further characterise these patterns.

\subsection*{Adjacent-stage differences and functional enrichment}
Study-adjusted models identified 857 substantial expression differences for F0--F1, 1,048 for F1--F2, 1,116 for F2--F3 and 781 for F3--F4 within the selected fibrosis-associated gene set (Fig.~\ref{fig:4}a). Substantial differences met both \(q<0.05\) and \(|\beta|\geq0.2\) within-study standard-deviation units. Increasing/decreasing counts were 822/35, 459/589, 792/324 and 723/58, respectively. F2--F3 had the largest number of substantial differences in the full analysis, while F3--F4 had the largest median absolute coefficient.

\clearpage
\begin{figure}[H]\centering
\begin{subfigure}[t]{0.98\linewidth}\centering
\projectimage{task3_transition_comparison.pdf}{0.72\textheight}\caption{}\end{subfigure}
\caption{\textbf{Adjacent-stage comparisons and functional enrichment.} (a) Full-data substantial-gene counts, median counts across 50 sample-size-matched refits, and median absolute study-adjusted coefficients for the four contrasts. Substantial changes require \(q<0.05\) and \(|\beta|\geq0.2\); matched refits use 30 patients per stage with fixed study composition. Matched-refit bars show medians without error bars. Coefficients are in within-study standard-deviation units. (b) Hallmark enrichment across six direction groups and five full clusters, using the 14,794-symbol background. Colour indicates \(-\log_{10}(q)\), capped at 10; blank/lightest cells have no significant term in the results. Confidence-core comparisons are reported separately.}\label{fig:4}
\end{figure}
\clearpage
\begin{figure}[H]\ContinuedFloat\centering
\setcounter{subfigure}{1}
\begin{subfigure}[t]{0.98\textwidth}\centering
\projectimage{task4_heatmap_hallmark.pdf}{0.79\textheight}\caption{}\end{subfigure}
\caption{\textbf{Adjacent-stage comparisons and functional enrichment (continued).} (b) Hallmark enrichment heatmap. Colour represents enrichment significance, with \(-\log_{10}(q)\) capped at 10.}
\end{figure}
\clearpage

The first three comparisons included 152/160, 160/152 and 152/112 patient-level observations across five studies. F3--F4 included 86/34 observations across four studies represented at both stages. Median absolute coefficients were 0.160, 0.186, 0.224 and 0.352, respectively. Matched refits with 30 patients per stage produced mean substantial-gene counts of 12.42, 10.22, 77.02 and 369.46, with SDs of 67.6, 34.4, 154.2 and 383.4. Corresponding medians were 0, 0, 1 and 232.5. The differences between full and matched results, together with broad refit distributions, show that transition-level counts were sensitive to patient composition and statistical power. These comparisons remained conditional on prior fibrosis-gene selection.

Functional enrichment identified distinct annotation profiles across the direction groups and expression clusters (Fig.~\ref{fig:4}b). Fibrosis-only increasing genes were enriched for T-cell activation (80 overlapping genes; fold enrichment 2.281; q = 8.19 \ensuremath{\times} 10\ensuremath{^{-9}}) and Hallmark IL6–JAK–STAT3 signalling (19 genes; 3.556; q = 2.14 \ensuremath{\times} 10\ensuremath{^{-5}}). Genes with positive associations on both axes were enriched for matrix biology, including collagen fibril organisation (GO Biological Process; 29 genes; fold enrichment 5.165; q = 1.40 \ensuremath{\times} 10\ensuremath{^{-11}}), collagen formation (Reactome; 32 genes; 4.151; q = 1.47 \ensuremath{\times} 10\ensuremath{^{-10}}) and extracellular matrix organisation (Reactome; 74 genes; 2.956; q = 3.80 \ensuremath{\times} 10\ensuremath{^{-15}}). Among expression clusters, C2 was enriched for the Hallmark epithelial–mesenchymal transition set (74 genes; 4.215; q = 5.25 \ensuremath{\times} 10\ensuremath{^{-27}}), while C3 was enriched for E2F targets (25 genes; 8.307; q = 1.24 \ensuremath{\times} 10\ensuremath{^{-14}}). The decreasing C5 cluster was enriched for bile-acid metabolism (37 genes; 4.986; q = 2.66 \ensuremath{\times} 10\ensuremath{^{-15}}) and fatty-acid metabolism (36 genes; 3.606; q = 1.42 \ensuremath{\times} 10\ensuremath{^{-10}}). These results describe annotation over-representation and do not establish activation or suppression of the named processes.

Enrichment also varied with cluster-membership confidence. C1 and C2 had 370 and 518 significant GO Biological Process terms and 18 and 20 Hallmark terms, respectively. C4 had no significant GO Biological Process terms, but one KEGG, one Reactome and three Hallmark terms. Its confidence core had one significant Reactome term and no significant terms in the other collections. The full C4 Reactome result was intrinsic fibrin-clot formation (\(q=0.0114\)), whereas its core result was BMP signalling (\(q=0.00150\)); the two sets shared no significant Reactome term. In contrast, 19 of the 21 C2 core Hallmark terms also appeared in the full cluster, and all five C5 core Hallmark terms were shared with the full set. These differences supported retaining membership uncertainty alongside functional context in gene prioritisation.

\subsection*{Gene prioritisation, robustness and sensitivity}
The multi-evidence scoring framework (P, T, X, F and B components; see Methods) identified 447 Tier 1 genes at a score threshold of seven: 93 in C1, 262 in C2, seven in C3, three in C4 and 82 in C5. Selecting up to five qualifying genes per cluster produced a 23-gene shortlist (Table~\ref{tab:3}). C4 contributed only three candidates because only three of its genes met the Tier 1 threshold. PDGFRA, CACNA1C, TPM4, PLK3 and TAX1BP3 scored nine; the remaining shortlisted genes scored seven or eight. The shortlist represented all five expression patterns while retaining the operational evidence threshold.

\begingroup\fontsize{9}{11}\selectfont\setlength{\tabcolsep}{3pt}
\begin{longtable}{>{\raggedright\arraybackslash}p{0.11\textwidth}>{\raggedright\arraybackslash}p{0.06\textwidth}>{\raggedright\arraybackslash}p{0.08\textwidth}>{\raggedright\arraybackslash}p{0.13\textwidth}>{\raggedright\arraybackslash}p{0.13\textwidth}>{\raggedright\arraybackslash}p{0.14\textwidth}>{\raggedright\arraybackslash}p{0.16\textwidth}}
\caption{Final 23-gene shortlist}\label{tab:3}\\
\toprule
\textbf{Gene} & \textbf{Cluster} & \textbf{Score /9} & \textbf{Median fibrosis rho} & \textbf{Membership confidence} & \textbf{Valid bootstrap denominator} & \textbf{Largest absolute contrast} \\
\midrule\endfirsthead
\toprule
\textbf{Gene} & \textbf{Cluster} & \textbf{Score /9} & \textbf{Median fibrosis rho} & \textbf{Membership confidence} & \textbf{Valid bootstrap denominator} & \textbf{Largest absolute contrast} \\
\midrule\endhead
\bottomrule\endfoot
IL32 & C1 & 8 & 0.522 & 0.83 & 100 & F0\ensuremath{\rightarrow}F1 \\
EPHA3 & C1 & 8 & 0.498 & 0.46 & 100 & F2\ensuremath{\rightarrow}F3 \\
LYZ & C1 & 8 & 0.487 & 0.83 & 100 & F0\ensuremath{\rightarrow}F1 \\
NFASC & C1 & 8 & 0.532 & 0.54 & 100 & F2\ensuremath{\rightarrow}F3 \\
PLPP4 & C1 & 8 & 0.533 & 0.58 & 100 & F2\ensuremath{\rightarrow}F3 \\
PDGFRA & C2 & 9 & 0.417 & 0.68 & 97 & F2\ensuremath{\rightarrow}F3 \\
CACNA1C & C2 & 9 & 0.447 & 0.89 & 97 & F3\ensuremath{\rightarrow}F4 \\
TPM4 & C2 & 9 & 0.403 & 0.89 & 97 & F3\ensuremath{\rightarrow}F4 \\
PLK3 & C2 & 9 & 0.414 & 0.89 & 97 & F3\ensuremath{\rightarrow}F4 \\
TAX1BP3 & C2 & 9 & 0.408 & 0.92 & 97 & F3\ensuremath{\rightarrow}F4 \\
TREM2 & C3 & 8 & 0.469 & 0.67 & 99 & F0\ensuremath{\rightarrow}F1 \\
CDC45 & C3 & 7 & 0.401 & 0.75 & 99 & F0\ensuremath{\rightarrow}F1 \\
HTR7 & C3 & 7 & 0.44 & 0.71 & 99 & F2\ensuremath{\rightarrow}F3 \\
FABP5 & C3 & 7 & 0.464 & 0.73 & 99 & F0\ensuremath{\rightarrow}F1 \\
TYMS & C3 & 7 & 0.489 & 0.61 & 99 & F0\ensuremath{\rightarrow}F1 \\
OAT & C4 & 7 & -0.398 & 0.38 & 100 & F2\ensuremath{\rightarrow}F3 \\
KNG1 & C4 & 7 & -0.312 & 0.71 & 100 & F1\ensuremath{\rightarrow}F2 \\
F2 & C4 & 7 & -0.307 & 0.56 & 100 & F1\ensuremath{\rightarrow}F2 \\
SULT1E1 & C5 & 8 & -0.414 & 0.99 & 96 & F3\ensuremath{\rightarrow}F4 \\
SLC38A4 & C5 & 8 & -0.43 & 0.99 & 96 & F3\ensuremath{\rightarrow}F4 \\
CYP2C19 & C5 & 8 & -0.443 & 0.76 & 96 & F1\ensuremath{\rightarrow}F2 \\
ACADSB & C5 & 8 & -0.447 & 0.89 & 96 & F1\ensuremath{\rightarrow}F2 \\
MME & C5 & 8 & -0.421 & 0.99 & 96 & F3\ensuremath{\rightarrow}F4 \\
\end{longtable}\endgroup

The largest absolute contrast in Table~\ref{tab:3} is the gene-level study-adjusted coefficient and need not coincide with the transition defining its cluster's average pattern. Cluster-membership confidence ranged from 0.38 for OAT to 0.99 for several C5 candidates. EPHA3, NFASC, PDGFRA and TREM2 had confidence values of 0.46, 0.54, 0.68 and 0.67, respectively. High prioritisation scores therefore did not ensure equally strong support for individual cluster assignments or stage timing.

LOSO analysis retained statistical significance for 442 of 447 Tier 1 genes in all five refits, whereas 247 retained an absolute median correlation of at least 0.30 in every refit. Within the 23-gene shortlist, all candidates remained significant in all five refits, and 18 retained the correlation threshold throughout. These findings distinguish significance retention from effect-magnitude stability and constitute internal robustness evidence for previously selected genes.

Prioritisation was also sensitive to the examined label and normalisation choices. Sixteen genes were shortlisted in all four specified variants. Under the final combined processing, removing the F feature while retaining the Tier 1 threshold of seven yielded 286 Tier 1 genes and a 21-gene shortlist, compared with 447 and 23 under the full score. This stricter feature-removal comparison reduced the maximum possible score from nine to eight. The alternative outputs are sensitivity evidence; the combined processing with the complete P/T/X/F/B score remained the final analysis. The prioritised genes were subsequently examined in their curated graph context.

\subsection*{Knowledge-graph annotation of prioritised genes}
The focused annotation subgraph contained 42 Tier 1 genes, 12 selected Pathway/BioProcess annotations and 154 gene–annotation connections (Fig.~\ref{fig:5}). Enriched annotations included extracellular matrix organisation (27 Tier 1 genes; fold enrichment 7.77; q = 1.06 \ensuremath{\times} 10\ensuremath{^{-13}}), cell–cell adhesion (32 genes; 5.35; q = 8.60 \ensuremath{\times} 10\ensuremath{^{-12}}), angiogenesis (36 genes; 4.69; q = 8.60 \ensuremath{\times} 10\ensuremath{^{-12}}), integrin cell surface interactions (20 genes; 8.38; q = 1.08 \ensuremath{\times} 10\ensuremath{^{-10}}) and signalling by PDGF (12 genes; 12.81; q = 1.28 \ensuremath{\times} 10\ensuremath{^{-8}}). The selected set also contained the cross-tissue annotation "Developmental Lineage of Pancreatic Ductal Cells", which does not by itself establish that process in the liver samples.

\begin{figure}[H]\centering
\includegraphics[width=\linewidth,height=0.80\textheight,keepaspectratio]{12_mechanistic_subgraph.pdf}
\caption{\textbf{Selected Tier 1 gene--annotation connections.} The filtered network contains 42 genes, 12 selected Pathway/BioProcess annotations and 154 links. Genes reach at least three selected annotations or are shortlisted genes reaching at least one. PDGFRA reaches five annotations, while EPHA3, NFASC and TREM2 each reach one; these are the four represented shortlist genes. Annotations were selected by over-representation and redundancy filtering as described in Methods. Gene nodes flank the annotation column; node positions are chosen for readability. Connections denote curated annotations rather than physical interactions or a demonstrated causal cascade.}\label{fig:5}
\end{figure}

Only four of the 23 shortlisted genes appeared in this filtered view. PDGFRA reached five selected annotations; EPHA3, NFASC and TREM2 each reached one. The other displayed genes came from the wider Tier 1 set and met the rule of reaching at least three selected annotations. The network therefore illustrated annotation sharing within a selected Tier 1 subset rather than convergence across the complete shortlist. These annotation links provide biological context for the prioritised genes and do not establish causal mechanisms.




\section*{Discussion}
\subsection*{Cross-study associations and their contribution to fibrosis research}
The principal finding is that fibrosis severity is associated with several distinct transcriptional patterns, rather than a single uniform expression response. Integrating eight liver transcriptomic datasets with biomedical annotations identified 3,845 fibrosis-associated genes and organised their F0--F4 profiles into five patterns. The resulting resource links association direction, stage-specific expression differences, membership uncertainty and biological annotations for candidate-gene investigation. Its contribution is the joint organisation of these evidence types within a stage-informed graph, rather than a demonstrated improvement in clinical prediction.

This analysis complements previous work relating liver gene expression to histological severity and molecular heterogeneity \cite{ref4}. Here, cross-study directional agreement was considered alongside the full F0--F4 expression profile and annotation context. The distinction between fibrosis and ordered disease-group associations was also retained: 1,822 genes were shared, including 30 with discordant directions. Because the two analyses involved different study sets and clinical labels, a fibrosis-only designation describes an operational selection outcome, not proven biological specificity. These distinctions provide the basis for interpreting the expression patterns and their functional context without collapsing them into one progression signature.

\subsection*{Biological interpretation of the expression patterns}
The increasing patterns were accompanied by extracellular-matrix, collagen, adhesion and immune-related annotations, whereas the late-decreasing C5 pattern was enriched for bile-acid and fatty-acid metabolism gene sets. Together, these findings suggest that higher fibrosis stages are associated with differences in both tissue-remodelling and metabolic expression programmes. Enrichment identifies over-represented annotations; it does not establish pathway activity, epithelial-to-mesenchymal conversion or the direction of a causal biological effect. Bulk-liver measurements also cannot separate changes within cells from changes in the proportions of the cells sampled.

Independent single-cell research identified scar-associated TREM2-positive macrophages and PDGFR\ensuremath{\alpha}-positive mesenchymal populations in human cirrhotic liver \cite{ref7}. The inclusion of TREM2 and PDGFRA in the present shortlist is compatible with that cellular context and motivates cell-resolved follow-up. It does not demonstrate the cellular origin or functional role of their associations in these cohorts. The pancreatic ductal developmental annotation likewise indicates shared curated gene membership rather than evidence of that developmental process in liver tissue. Such biological interpretations therefore need to be weighed against the stability of the underlying expression patterns.

\subsection*{Expression-pattern stability and adjacent-stage evidence}
The five-cluster solution was stable under gene subsampling, but less stable when patients were resampled: mean ARI values at \(k=5\) were 0.983 and 0.592, respectively. Only 1,838 of the 3,845 genes met the conditional membership-confidence threshold. These findings indicate that reproducible average shapes can coexist with uncertain assignments for individual genes. The highest silhouette score occurred at \(k=3\), so the five-pattern representation should be understood as the outcome of the stated stability rule rather than a uniquely established biological partition. Reporting membership confidence alongside cluster labels is consequently important when selecting genes for follow-up.

Adjacent-stage contrasts add resolution to these average shapes, but their apparent intensity depends on sample size and patient composition. Although F3--F4 had the largest median absolute coefficient, only 34 F4 observations contributed to that comparison, and sample-size-matched refits produced widely varying substantial-gene counts. The contrast with the full-data counts limits interpretation of any single stage boundary as a biological switch. Study adjustment reduces one source of heterogeneity, while clinical confounding and cell-composition differences remain possible. Moreover, these cross-sectional comparisons describe differences between patients; they do not establish the sequence of molecular events within an individual. This distinction also defines what the prioritisation framework can support.

\subsection*{Candidate prioritisation and knowledge-graph context}
The multi-evidence score identified 447 Tier 1 genes, with a 23-gene shortlist representing all five expression patterns. Keeping the qualification threshold fixed allowed the three qualifying C4 genes to represent that pattern without filling an arbitrary quota. Leave-one-study-out refits retained significance for 442 Tier 1 genes, but only 247 retained the effect-magnitude criterion throughout. All 23 shortlisted genes remained significant, whereas 18 retained the magnitude threshold. Thus, statistical-significance retention and effect stability provide different information, and neither alone establishes candidate validity in an external population.

The score components also share expression and annotation evidence, so a high total is not a set of independent confirmations. Sixteen genes appeared in all four label/transformation variants; removing the F component reduced the Tier 1 set to 286 genes under a lower maximum achievable score. These comparisons describe sensitivity to analytical choices rather than performance against an external benchmark. The focused graph adds interpretable annotation links, but only four shortlisted genes appear among its 42 displayed genes. PDGFRA reaches five selected annotations, while EPHA3, NFASC and TREM2 each reach one. Connectivity therefore supports a contextual view of a selected subset, not convergence of the entire shortlist. No comparison under identical datasets and evaluation criteria established superiority over another graph or prioritisation method. These boundaries inform the remaining limitations and the validation needed before broader use.

\subsection*{Study limitations and reproducibility}
The datasets were not selected through a documented systematic search, which limits claims of comprehensive coverage. Five ambiguous GSE130970 observations were excluded from disease-group analysis, but the remaining 73 classifications were derived from histology rather than recovered original diagnoses. Incomplete metabolic and alcohol information also prevents direct reconstruction of contemporary MASLD/MASH criteria \cite{ref3}. Demographic summaries cover only 327 patient units, and participant overlap across accessions remains unresolved. Repeat-biopsy handling reduces duplicate representation within studies but does not establish independence between cohorts. Unequal stage representation and incomplete covariate information further limit causal interpretation of the observed associations.

Reproducibility is supported by explicit input paths, retained-file checksums, parameters and analysis outputs. However, the source manifest was generated after the analysis; it identifies retained file contents rather than providing a contemporaneous record of every acquisition and analytical decision. File modification dates do not establish official HGNC or OptimusKG release dates. The shared expression-analysis core also does not replace a complete archive of all graph measurements. Candidate selection was not fully nested within the robustness assessment, and independent-cohort, longitudinal and functional validation remain absent. The findings therefore support prioritised hypotheses and an organised research resource, with subsequent validation required to establish their generalisability and biological roles.

\section*{Conclusion}
Liver fibrosis severity was associated with heterogeneous expression patterns that can be organised alongside cross-study consistency, stage contrasts and curated biological annotations. The practical value of this framework is that candidate genes can be examined together with the evidence supporting their ranking and the uncertainty surrounding their stage-pattern assignment. Association strength, cluster confidence and graph connectivity contribute different forms of information and should remain distinguishable when the resource is used for further research. The present evidence supports gene prioritisation and biological-context assessment, rather than validated patient prediction or causal conclusions.

Three priorities follow from these boundaries. First, independent cohorts with harmonised clinical labels and longitudinal sampling should test the reproducibility and temporal interpretation of the candidate associations, using nested selection where performance is evaluated. Second, cell-resolved or spatial measurements should determine whether the bulk-liver patterns reflect cellular composition, cell-state changes or both, followed by targeted functional experiments. Third, an immutable public release of the analysis code, graph exports, source manifests and complete supplementary outputs should support external reproduction and fair comparisons under shared evaluation conditions. These steps would establish the scope in which the resource and its prioritised genes can be used reliably.

\section*{Materials and Methods}
\subsection*{Study design and data sources}
We conducted a retrospective computational study integrating public human liver transcriptomic data with biomedical knowledge to identify genes associated with fibrosis severity, characterise expression patterns across fibrosis stages F0–F4, and prioritise candidates for further biological investigation. The overall workflow is summarised in Fig.~\ref{fig:6}. Eight datasets deposited in the NCBI Gene Expression Omnibus (GEO) were analysed \cite{ref8}: GSE126848 (57 samples), GSE130970 (78), GSE135251 (216), GSE162694 (143), GSE167523 (98), GSE193066 (164), GSE240729 (67) and GSE269412 (262). Together, these datasets comprised 1,085 samples and 1,027 recorded patient identifiers. Patient identifiers were interpreted within each study; participant independence across accessions was not established.

Deposited supplementary expression matrices and corresponding series-matrix metadata were used as the primary inputs. These included gene-count matrices, a transcripts-per-million (TPM) matrix, per-sample expression files and a GCT-format matrix, which required study-specific preprocessing. Available metadata included fibrosis stage, source disease labels, age, sex, patient identity and biopsy order, with coverage varying between datasets. No new sequencing, read alignment or transcript quantification was performed. Sample availability determined the study size, and eligibility for each analysis depended on the required clinical annotations. Dataset characteristics, deposited expression filenames and annotation coverage are summarised in Table~\ref{tab:1}. These differences in metadata availability informed the clinical harmonisation and patient-handling procedures described below. The datasets were not identified through a documented systematic search protocol. They comprised human liver tissue studies with publicly available sample-level expression files and metadata permitting linkage of measurements to samples; relevant fibrosis-stage or disease-group annotations then determined eligibility for each clinical analysis, rather than being a requirement met by every dataset.

\begingroup\fontsize{9}{11}\selectfont\setlength{\tabcolsep}{3pt}
\begin{longtable}{>{\raggedright\arraybackslash}p{0.1\textwidth}>{\raggedright\arraybackslash}p{0.075\textwidth}>{\raggedright\arraybackslash}p{0.09\textwidth}>{\raggedright\arraybackslash}p{0.075\textwidth}>{\raggedright\arraybackslash}p{0.12\textwidth}>{\raggedright\arraybackslash}p{0.37\textwidth}}
\caption{Public liver transcriptomic datasets and clinical annotation coverage}\label{tab:1}\\
\toprule
\textbf{Accession} & \textbf{Samples} & \textbf{Recorded patient IDs} & \textbf{Staged biopsies} & \textbf{Samples with disease labels} & \textbf{Deposited expression file} \\
\midrule\endfirsthead
\toprule
\textbf{Accession} & \textbf{Samples} & \textbf{Recorded patient IDs} & \textbf{Staged biopsies} & \textbf{Samples with disease labels} & \textbf{Deposited expression file} \\
\midrule\endhead
\bottomrule\endfoot
GSE126848 & 57 & 57 & 0 & 57 & \nolinkurl{GSE126848_Gene_counts_raw.txt.gz} \\
GSE130970 & 78 & 78 & 78 & 73 & \nolinkurl{GSE130970_all_sample_salmon_tximport_TPM_entrez_gene_ID.csv.gz} \\
GSE135251 & 216 & 216 & 216 & 216 & \nolinkurl{GSE135251_RAW.tar} \\
GSE162694 & 143 & 143 & 143 & 143 & \nolinkurl{GSE162694_raw_counts.csv.gz} \\
GSE167523 & 98 & 98 & 0 & 98 & \nolinkurl{GSE167523_Raw_gene_counts_matrix.txt.gz} \\
GSE193066 & 164 & 106 & 164 & 0 & \nolinkurl{GSE193066_NAFLD.HUn164.gct.gz} \\
GSE240729 & 67 & 67 & 67 & 0 & \nolinkurl{GSE240729_merged.counts.67samp.txt.gz} \\
GSE269412 & 262 & 262 & 0 & 0 & \nolinkurl{GSE269412_DRX0FLO_rnaseq_count_mat.txt.gz} \\
\textbf{Total} & \textbf{1,085} & \textbf{1,027} & \textbf{668} & \textbf{587} & --- \\
\end{longtable}\endgroup

Counts describe the full resource before patient-level repeat-biopsy resolution. Recorded patient IDs are counted within studies and do not establish distinct participants across accessions. Samples with disease labels include broad NAFLD and obese categories that are excluded from the ordered disease-group analysis. Five GSE130970 cases are unclassified; the remaining 73 disease labels are histology-derived. Staged and labelled sample counts therefore differ from the analysis populations. For the final ordered disease-group analysis, GSE126848 contributed 45 eligible observations after exclusion of 12 obese-only samples; GSE130970 contributed 73 after the five flagged classifications were set to unclassified; and GSE167523 contributed all 98. Thus, the final total was 45 + 73 + 98 = 216 observations.

\subsection*{Clinical annotation and patient handling}
Expression sample identifiers were linked to GEO sample records using study-specific identifier, title and description fields. Available clinical information was harmonised into a common structure containing study accession, sample identifier, patient identifier, biopsy order, fibrosis stage and disease group. Original clinical fields were retained alongside harmonised values and annotation provenance. Fibrosis severity was represented by ordinal stages F0--F4. Disease groups were ordered as control = 0, nonalcoholic fatty liver (NAFL) = 1 and nonalcoholic steatohepatitis (NASH) = 2. Broad NAFLD labels without an informative subtype, obese-only labels and unclassified records were excluded from the ordered disease-group analysis. Original diagnoses were preserved without assuming equivalence to contemporary MASLD/MASH criteria \cite{ref3}.

For GSE130970, individual author-assigned diagnoses were not recovered; disease-group labels were therefore derived from deposited histological component scores. Steatosis grade zero defined a control proxy, positive steatosis with positive hepatocyte ballooning defined a NASH proxy, and positive steatosis with zero ballooning defined a NAFL proxy. These operational labels did not reproduce a complete histopathological diagnostic assessment \cite{ref4}. Five flagged samples (GSM3758028, GSM3758038, GSM3758039, GSM3758073 and GSM3758074) were assigned unclassified disease-group status while retaining their original histological information and fibrosis stages. The remaining 73 labels were retained as histology-derived classifications. This unclassified status describes the final primary configuration; in the sensitivity analyses described below, the baseline and normalisation-only configurations retained the original classifications of these five samples. GSE193066 contributed fibrosis-stage information but had no usable disease-group diagnoses.

A common patient-level policy was applied to repeated biopsies. When all recorded fibrosis stages for a patient were known and identical, standardised gene-expression measurements were averaged across biopsies. When stages differed, only the first biopsy was retained; unknown stages were not assumed to be equal and also led to first-biopsy selection. Biopsy selection preceded filtering for individual stage comparisons, preventing a later biopsy from entering a comparison solely because the first biopsy fell outside the stages being compared. This policy was adopted so that patients with repeated biopsies did not receive disproportionate weight. GSE193066 contained 164 biopsies from 106 recorded patients, including 58 repeat-biopsy pairs \cite{ref10}: 28 pairs had matching stages and were averaged, whereas 30 had differing stages and contributed only their first biopsy.

The same patient-level policy was used for association discovery, stage-expression summaries, clustering and resampling, adjacent-stage comparisons and leave-one-study-out analyses. Each recorded patient contributed one analytical observation within a study, while all original biopsies and raw measurements were retained in the underlying resource. Study-level eligibility for association analysis required at least 25 patient-level observations and variation in the relevant clinical axis. The resulting analysis populations comprised 610 fibrosis-staged observations across GSE130970, GSE135251, GSE162694, GSE193066 and GSE240729, and 216 observations with ordered disease-group labels across GSE126848, GSE130970 and GSE167523. GSE135251 and GSE162694 contained only control and broad NAFLD samples, leaving no variation on the ordered disease-group axis.

\subsection*{Gene harmonisation and expression preprocessing}
Matrix, GCT and per-sample archive inputs were parsed into sample--gene measurements. Missing and non-numeric measurements were removed, the first occurrence of duplicate rows sharing a raw gene identifier was retained, and measured zeros were preserved. The HGNC complete set was used to map approved symbols, Entrez identifiers and unambiguous previous symbols or aliases to canonical Ensembl gene identifiers. Ensembl version suffixes and \texttt{\_PAR\_Y} rows were removed. Ensembl identifiers lacking HGNC annotation were retained without adding unsupported annotations, and mapping failures were recorded. Measurements from distinct raw identifiers mapping to the same canonical gene were summed. The graph retained the union of measured genes, whereas association discovery was restricted to the 14,794 genes measured in all eight studies, establishing a common eligibility universe.

The GSE130970 TPM matrix was transformed as \(\log_2(\mathrm{TPM}+1)\). Other deposited count matrices were transformed as \(\log_2(\mathrm{CPM}+1)\), with sample totals calculated before gene mapping and filtering. GSE193066 was treated separately because its deposited GCT contained linear DESeq2 relative-log-expression (RLE)-normalised counts. For this study, mapped aliases were summed before applying \(\log_2(\mathrm{count}+1)\), preserving the deposited normalisation without additional CPM rescaling. DESeq2 is cited for the source normalisation context \cite{ref9}; no DESeq2 differential-expression model was fitted in this analysis.

Transformed measurements were standardised separately for each gene within each study by subtracting the sample mean and dividing by the sample standard deviation (\(\mathrm{ddof}=1\)). Standardised expression was set to zero where the standard deviation was undefined. For GSE193066, standardisation used all 164 source biopsies before the patient-level biopsy policy was applied. This scaling provided a common expression unit for downstream summaries while preserving the study structure used in association testing. Standardised expression (\texttt{value\_z}) and derived effects are in within-study standard-deviation units rather than fold changes. No fitted cross-study batch-correction procedure was applied. The processing order was raw-identifier duplicate resolution (first occurrence retained), identifier mapping and exclusion of unmappable or PAR-Y records, summation of distinct raw identifiers mapping to the same sample--Ensembl gene, study-appropriate transformation, within-study gene standardisation, and finally patient-level biopsy selection and aggregation for analysis. For CPM-based transformations, library totals were calculated from the parsed measurements before mapping and gene filtering; the transformation itself followed mapped-gene summation. Equal-stage repeat biopsies were averaged on the standardised-expression scale, without subsequent patient-level restandardisation.

\subsection*{Cross-study association analysis}
Within each eligible study, Spearman rank correlation related patient-level standardised expression to fibrosis stage or ordered disease-group labels, with average ranks used for ties. Rank correlation was chosen because both clinical axes are ordinal and expression changes need not be linear. Two-sided p-values were calculated using the t approximation,
\[
t=\rho\sqrt{\frac{n-2}{1-\rho^2}},
\]
with \(n-2\) degrees of freedom. Correlations were clipped to \(\pm0.9999999\) for numerical calculation. A within-study association was testable only when expression measurements were available for all included patients and the ranked expression vector had nonzero variance. Genes required testable associations in at least three eligible studies. For each gene, the median study correlation summarised association strength and direction, and Fisher's statistic combined the contributing p-values:
\[
X_g^2=-2\sum_{d=1}^{D_g}\log(p_{gd}),\qquad X_g^2\sim\chi^2_{2D_g},
\]
where \(D_g\) is the number of testable studies for gene \(g\). P-values were clipped below at \(10^{-300}\) before combination. Fisher's combined p-value used the available two-sided within-study p-values without incorporating their correlation signs or the median correlation. Directional consistency was evaluated separately: each testable study's correlation was required to have the same positive/non-positive class as the median correlation. BH adjustment preceded application of this directional selection criterion.

Benjamini--Hochberg (BH) adjustment \cite{ref11} was applied separately to the fibrosis and disease-group test families, comprising 14,794 and 14,784 genes, respectively. The common measured-gene universe and the axis-specific testing family were distinct: ten common genes had constant expression among the 45 eligible GSE126848 observations and therefore had undefined Spearman correlations in that study. With only three eligible disease-group studies, these genes failed the minimum-three-testable-studies rule and were excluded from the family entering BH adjustment. Selection required \(q<0.05\) and agreement with the median association direction in every testable study. No minimum absolute-correlation threshold was imposed during discovery. Genes were classified as fibrosis-only, disease-group-only or shared according to membership in the two selected sets; shared genes were further distinguished by concordant or discordant directions. These labels describe statistical selection across different analysis populations rather than biological specificity. Exported stage means averaged patient observations, and median correlations were rounded to three decimal places in downstream summary and scoring inputs.

\subsection*{Fibrosis-stage expression profiles and clustering}
For each selected fibrosis-associated gene, mean standardised expression was calculated at each stage F0--F4 using the patient-level observations. The resulting five-stage profiles were row-standardised (\(\mathrm{ddof}=0\)) before K-means clustering, so clustering emphasised profile shape rather than absolute expression magnitude. Final fits used 50 initialisations and random seed 42. Candidate values \(k=3\)--8 were assessed using silhouette scores, gene-subset stability and patient-bootstrap stability. For each candidate k, 50 refits used random subsets containing 80\% of genes, and a separate 50 refits resampled patients within studies with replacement. Diagnostic fits used 10 initialisations. Stability was measured by the adjusted Rand index (ARI). The implemented selection rule retained the largest k with mean gene-subset \(\mathrm{ARI}\geq0.95\), yielding \(k=5\). Patient-bootstrap ARI and silhouette scores were reported separately to distinguish robustness to gene selection from sensitivity to patient composition. Supplementary Table~\ref{tab:S1} reports the numerical diagnostics for every evaluated k, with gene-subset and patient-bootstrap means and standard deviations presented separately.

An additional 100 patient bootstraps assessed centroid-shape reproducibility and individual gene membership. Hungarian assignment matched bootstrap centroids to reference centroids using their correlations. A gene was assessable in a bootstrap only when the matched centroid for its reference cluster had correlation \(r\geq0.5\). Membership confidence was the fraction of assessable bootstraps in which the gene returned to its matched reference cluster; the valid denominator was retained for each gene. Confidence-core sets required membership confidence \(\geq0.8\). As a further descriptive sensitivity check, stage profiles estimated by ordinary least squares with categorical stage and study terms were row-standardised, reclustered and compared with the main profiles using ARI and per-gene profile correlations. Cluster names were assigned from the fitted expression patterns. All profiles describe cross-sectional stage associations rather than longitudinal trajectories within patients.

\subsection*{Adjacent-stage comparisons and matched refits}
Adjacent-stage comparisons were performed for F0--F1, F1--F2, F2--F3 and F3--F4 within the selected fibrosis-associated gene set. For each contrast, only studies containing patients at both stages contributed. Standardised expression was modelled by ordinary least squares, with study-specific intercepts accounting for between-study differences in baseline expression:
\[
z_{ig}=\alpha_g+\beta_g H_i+\sum_{d=2}^{D}\gamma_{gd}I(\mathrm{study}_i=d)+\varepsilon_{ig},
\]
where \(H_i\) indicates membership in the higher stage and \(\beta_g\) estimates the study-adjusted expression difference. Two-sided t tests evaluated the stage coefficient, and BH adjustment was applied across the selected fibrosis gene set within each contrast. A substantial change required \(q<0.05\) and \(|\beta_g|\geq0.2\) within-study standard-deviation units. Study adjustment did not comprehensively control for demographic, clinical or cell-composition differences.

Matched refits assessed sensitivity to unequal stage sizes. Four studies were represented at every stage; their minimum patient counts across stages determined a fixed allocation of two patients per stage from GSE130970, 14 from GSE135251, eight from GSE162694 and six from GSE240729, giving 30 patients per stage. Fifty draws sampled without replacement within each study--stage pool using seed 42, and the same models were refitted. For each gene and transition, effect reproducibility was the fraction of draws with \(|\beta|\geq0.2\) and the same sign as the full-data coefficient. Matched-draw statistical significance was recorded separately. The matched design compared contrasts at common sample sizes, but all comparisons remained conditional on prior selection for fibrosis association.

\subsection*{Functional enrichment analysis}
Over-representation analysis used local MSigDB 2024.1.Hs human-symbol gene sets for GO Biological Process, KEGG legacy, Reactome and Hallmark \cite{ref12}. Tests were performed with GSEApy \cite{ref13} using a one-sided hypergeometric distribution. The background comprised the 14,794 eligible unique gene symbols corresponding one-to-one to the 14,794 Ensembl identifiers in the common eligibility universe. Both namespaces were taken from the retained harmonised gene records; no additional symbol-to-Ensembl conversion was performed during enrichment testing. Gene sets were intersected with this background and retained when the resulting set size was 10--500; input groups required at least five genes. BH correction was applied separately within each group and annotation collection, with \(q<0.05\) defining significant enrichment. This tested whether annotated genes were over-represented relative to the measured eligibility universe, rather than testing coordinated changes in a ranked whole-transcriptome list.

Tests covered six concordant direction groups---increasing and decreasing fibrosis-only, disease-group-only and shared genes---together with the five full expression clusters and their five confidence-core sets. Shared genes with discordant association directions were excluded from the shared increasing/decreasing groups. Enrichment outputs retained overlap counts, gene-set sizes, fold enrichment, adjusted p-values and contributing genes. Full-cluster and core results were compared to assess sensitivity to uncertain cluster membership. Groups below the minimum size were labelled untested, distinguishing them from tested groups with no significant terms. Enrichment and overlapping annotation membership were interpreted as biological context rather than demonstrated pathway activation.

\subsection*{Gene prioritisation and robustness assessment}
Fibrosis-associated genes were prioritised using the operational score
\[
S_g=P_g+T_g+X_g+F_g+B_g,
\]
with a maximum of nine points. Scoring used the exported median correlations and stage means, recorded to three decimal places; the endpoint difference was also rounded to three decimal places before threshold comparison. All stated lower thresholds were inclusive. \(P_g\) represented association magnitude: two points for \(|\operatorname{median}\rho|\geq0.40\), one for \(0.30\leq|\operatorname{median}\rho|<0.40\), and zero otherwise. Effect magnitude was used to discriminate genes already passing the discovery significance criterion. \(T_g\) represented the endpoint expression difference: two points for \(|\overline{z}_{F4}-\overline{z}_{F0}|\geq1\), one for a difference of at least 0.5 but less than 1 standard-deviation unit, and zero otherwise. The endpoint direction agreed when the rounded F4-minus-F0 difference and the exported median fibrosis correlation were both positive or both non-positive. Disagreement assigned zero before magnitude thresholds were evaluated; a zero endpoint difference also received zero because its magnitude was below 0.5. Cluster direction did not override this rule. This feature quantified endpoint separation, while \(X_g\) captured evidence at the largest absolute adjacent-stage coefficient, including intermediate-stage differences.

For \(X_g\), the selected comparison was the one with the largest absolute available full-data study-adjusted coefficient. Exact ties retained the first comparison in the tie-breaking order F0--F1, F1--F2, F2--F3, F3--F4. Two points required both a substantial full-data difference (\(q<0.05\) and \(|\beta|\geq0.2\)) and same-sign effects with \(|\beta|\geq0.2\) in at least 50\% of matched draws; one point required exactly one condition, and zero required neither. Same-sign reproducibility referred to the selected full-data coefficient, not to the median fibrosis correlation. Matched-draw significance was not scored. Implementation checks are described in the Supplementary Methods. \(F_g\) awarded one point for membership in the fibrosis-only selected set and zero otherwise. \(B_g\) counted annotation collections in which the gene symbol appeared in the overlapping-gene list of at least one significantly enriched term (BH \(q<0.05\)) for either its own concordant direction group or its full expression cluster. All significant terms were eligible, rather than only terms displayed among the top results; confidence-core enrichment did not contribute. Each of GO Biological Process, KEGG legacy, Reactome and Hallmark counted at most once, irrespective of the number of terms or whether both routes supported membership. Discordant shared genes had no concordant direction-group route but could receive support through their full cluster. Two or more collections awarded two points, one collection awarded one point, and no collection awarded zero. Tier 1 required \(S_g\geq7\), Tier 2 required \(5\leq S_g<7\), and Tier 3 comprised the remaining genes. Up to five Tier 1 genes per cluster were shortlisted by descending total score and then absolute median fibrosis correlation. A cluster with fewer than five qualifying genes contributed fewer candidates. Cluster quotas retained representation of the different expression patterns. These features share expression and annotation information, so the tiers are prioritisation categories rather than independent confirmations of target validity.

Leave-one-study-out (LOSO) analysis tested whether selected genes depended on any single study. It omitted each of the five fibrosis studies in turn, repeated the patient-level association calculations, and applied BH correction across the full 14,794-gene eligible family in every refit. For previously selected Tier 1 genes, summaries recorded retention of q < 0.05, association direction and |median \ensuremath{\rho}| \ensuremath{\geq} 0.30. Significance retention and effect-magnitude retention were assessed separately. Full gene selection and prioritisation were not repeated within an independent held-out prediction design; LOSO therefore assessed internal robustness.

Sensitivity analyses compared four label/transformation configurations. The baseline retained the original derived disease classifications for all 78 GSE130970 samples and applied CPM-style rescaling to GSE193066. The labels-only variant marked the five flagged disease classifications unclassified while retaining the baseline transformation; the normalisation-only variant preserved deposited GSE193066 RLE normalisation while retaining the baseline disease classifications. The combined variant applied both changes and was adopted as the final primary configuration. Consequently, those five observations did not contribute to the primary disease-group analysis, although their fibrosis stages remained available. Unless explicitly identified as sensitivity comparators, all reported analyses used this combined configuration. Biopsy handling, eligibility rules, \(k=5\), scoring thresholds, cluster quotas and annotation-selection rules were held constant. Additional analyses examined score-threshold changes, removal of each score component and an alternative cluster-top-score shortlist rule. A specific F-feature ablation set \(F_g=0\) while retaining the Tier 1 threshold of seven, reducing the maximum possible score to eight. These comparisons evaluated sensitivity to analytical choices rather than performance against external prediction benchmarks.

\subsection*{Knowledge-graph construction and analysis}
The expression layer and biomedical knowledge layer were integrated in Neo4j using Ensembl identifiers to match measured genes to OptimusKG entities \cite{ref5}. The OptimusKG source deposit was identified by DOI 10.7910/DVN/IYNGEV. The schema comprised Gene, Sample, Dataset, Disease, Drug, Pathway, BioProcess and Phenotype nodes. Sample-to-Gene \texttt{EXPRESSES} relationships retained raw and standardised measurements, and \texttt{IN\_DATASET} linked samples to their studies. Patient identity and biopsy order were Sample properties, allowing original biopsies to be retained while analytical patient units were defined separately.

The curated layer represented gene--disease associations, gene interactions, pathway and biological-process membership, disease hierarchies, disease--phenotype connections, and drug--gene and drug--disease relations. Gene--disease association scores originating from Open Targets through OptimusKG were retained at \(\mathrm{score}\geq0.1\). Symmetric gene interactions were deduplicated. Drug action labels and indication, contraindication and off-label relations retained their distinct meanings. Selected expression associations were stored separately as \texttt{TRACKS\_FIBROSIS} and \texttt{TRACKS\_INFLAMMATION}, retaining association direction, correlation, adjusted p-value and contributing-study count. The latter is an implementation name for disease-group associations rather than a direct inflammation measurement. 

\begingroup\raggedright
The final graph used 14 relationship types: \texttt{EXPRESSES}, \texttt{IN\_DATASET}, \texttt{ASSOCIATED\_WITH}, \texttt{INTERACTS\_WITH}, \texttt{INVOLVED\_IN}, \texttt{IN\_PATHWAY}, \texttt{ACTS\_ON}, \texttt{TREATS}, \texttt{CONTRAINDICATED\_IN}, \texttt{OFF\_LABEL\_FOR}, \texttt{HAS\_PHENOTYPE}, \texttt{PARENT\_OF}, \texttt{TRACKS\_FIBROSIS} and \texttt{TRACKS\_INFLAMMATION}.\par
\endgroup

 Verification checks compared derived properties and progression links with the analysis tables and confirmed preservation of raw expression measurements and curated relationship counts.

A focused annotation subgraph was constructed from Pathway and BioProcess nodes linked directly to Tier 1 genes. Hypergeometric tests evaluated Tier 1 over-representation against the same common measured-gene universe used for association eligibility, represented here by its 14,794 unique Ensembl identifiers. All 447 Tier 1 genes belonged to this background. Annotation sizes counted only members intersecting this universe, including genes without a given annotation in the background denominator. This Ensembl representation corresponded one-to-one with the 14,794-symbol functional-enrichment background. BH adjustment covered the 3,980 Pathway/BioProcess nodes reached by Tier 1 genes before size, connectivity and figure-eligibility filtering; 3,980 was the tested annotation family, not the graph-wide total number of annotation nodes. Eligible annotations contained 10--500 background genes, at least ten Tier 1 genes and \(q<0.05\). Annotations were ranked by adjusted p-value, and a candidate was skipped when its Tier 1 member set had Jaccard similarity greater than 0.5 to an already retained annotation. Twelve annotations were retained. Displayed genes reached at least three selected annotations or were shortlisted genes reaching at least one. The resulting subgraph contained 42 displayed Tier 1 genes and 154 gene--annotation relationships. These rules produced a compact view of annotated connections; annotation evidence used in prioritisation was not treated as independent validation.


\subsection*{Software and reproducibility}
The analyses used Python 3.12.14, pandas 2.3.3, NumPy 2.2.4, SciPy 1.16.3, scikit-learn 1.7.2, matplotlib 3.10.7 and GSEApy 1.3.1. The graph used Neo4j Enterprise 2026.09.0. Random seed 42 was used for clustering and matched refits as described above. Code, runtime requirements, result tables and provenance manifests document the computational procedures and inputs. The expression analysis input covers the common 14,794-gene core rather than the complete expression layer.

Expression-dataset acquisition dates and formal HGNC and OptimusKG release identifiers were not recorded. The OptimusKG source deposit is identified by DOI 10.7910/DVN/IYNGEV, and the functional-annotation collection is identified as MSigDB 2024.1.Hs. The retained reference files and their modification dates are described in the Supplementary Methods. File dates and checksums do not establish official release dates. No new human participant data were collected.




\section*{Data availability}
The source transcriptomic datasets are available from the Gene Expression Omnibus under accessions GSE126848, GSE130970, GSE135251, GSE162694, GSE167523, GSE193066, GSE240729 and GSE269412. The OptimusKG source deposit is identified by DOI \href{https://doi.org/10.7910/DVN/IYNGEV}{10.7910/DVN/IYNGEV}; reuse is subject to the terms of the contributing sources. The reference-file manifest, \nolinkurl{source_checksums.txt}, documents the retained HGNC and OptimusKG files. [AUTHOR\_CONFIRM: permanent repository links or DOIs for the final graph exports, processed analysis inputs, complete result tables and supplementary outputs; specify any access or redistribution restrictions.] The shared expression-analysis core covers the common measured-gene universe and is not a complete archive of all graph measurements.

\section*{Code availability}
[AUTHOR\_CONFIRM: public repository URL, immutable release or commit identifier, software licence and archival DOI for the code corresponding to the reported analyses.] The deposited release should include the runtime requirements and instructions linking the source inputs to the analysis outputs and figures.

\section*{Ethics approval and consent}
This study used publicly available human liver transcriptomic datasets and their associated annotations; no new participant recruitment or sample collection was undertaken. [AUTHOR\_CONFIRM: applicable institutional determination for secondary-data ethics review and the consent and data-use conditions governing the source datasets.]

\section*{AI-assisted tools}
Claude and Codex were used to assist with project support and manuscript preparation. [AUTHOR\_CONFIRM: specify the scope of AI-assisted analysis, coding and writing, the tools or versions where available, and the human verification performed. Confirm that the authors reviewed the final text and outputs and take responsibility for the manuscript.]

\section*{Acknowledgements}
[AUTHOR\_CONFIRM: contributors to acknowledge, the nature of their assistance and permission to name them.]

\section*{Funding}
[AUTHOR\_CONFIRM: funding organisations and grant identifiers, or a verified statement that no specific funding was received; describe the funders' role where applicable.]

\section*{Author contributions}
[AUTHOR\_CONFIRM: author initials and their contributions to conceptualisation, data curation, methodology, software, formal analysis, visualisation, supervision, writing and revision. Confirm approval of the submitted manuscript by all authors.]

\section*{Competing interests}
[AUTHOR\_CONFIRM: financial and non-financial competing-interest declarations for every author.]

\section*{Additional information}
Correspondence and requests for materials should be addressed to [CORRESPONDING AUTHOR]. Supplementary Table S1 reports cluster-number diagnostics; the Supplementary Methods describe implementation checks and reference-file provenance. [AUTHOR\_CONFIRM: permanent location and final inventory of the complete supplementary deposit.]

\begin{thebibliography}{99}
\bibitem{ref1} Angulo, P. et al. Liver Fibrosis, but No Other Histologic Features, Is Associated With Long-term Outcomes of Patients With Nonalcoholic Fatty Liver Disease. \textit{Gastroenterology} \textbf{149}, 389–397.e10 (2015). \url{https://doi.org/10.1053/j.gastro.2015.04.043}.
\bibitem{ref2} Sanyal, A. J. et al. Prospective Study of Outcomes in Adults with Nonalcoholic Fatty Liver Disease. \textit{N. Engl. J. Med.} \textbf{385}, 1559–1569 (2021). \url{https://doi.org/10.1056/NEJMoa2029349}.
\bibitem{ref3} Rinella, M. E. et al. A multisociety Delphi consensus statement on new fatty liver disease nomenclature. \textit{Hepatology} \textbf{78}, 1966–1986 (2023). \url{https://doi.org/10.1097/HEP.0000000000000520}.
\bibitem{singh2015} Singh, S. et al. Fibrosis progression in nonalcoholic fatty liver vs nonalcoholic steatohepatitis: a systematic review and meta-analysis of paired-biopsy studies. \textit{Clin. Gastroenterol. Hepatol.} \textbf{13}, 643--654.e1--9 (2015). \url{https://doi.org/10.1016/j.cgh.2014.04.014}.
\bibitem{govaere2020} Govaere, O. et al. Transcriptomic profiling across the nonalcoholic fatty liver disease spectrum reveals gene signatures for steatohepatitis and fibrosis. \textit{Sci. Transl. Med.} \textbf{12}, eaba4448 (2020). \url{https://doi.org/10.1126/scitranslmed.aba4448}.
\bibitem{piras2024} Piras, I. S. \& DiStefano, J. K. Comprehensive meta-analysis reveals distinct gene expression signatures of MASLD progression. \textit{Life Sci. Alliance} \textbf{7}, e202302517 (2024). \url{https://doi.org/10.26508/lsa.202302517}.
\bibitem{yang2026} Yang, W. et al. Integrative Transcriptomics Across Etiologies Reveals Common and Disease-Specific Fibrogenic Signatures in Liver Fibrosis. \textit{Can. J. Gastroenterol. Hepatol.} \textbf{2026}, e9254321 (2026). \url{https://doi.org/10.1155/cjgh/9254321}.
\bibitem{maiers2026} Maiers, J. L. et al. Conserved hepatic RNA signatures across multiple cohorts reveal novel mechanistic clues for advanced fibrosis in human MASLD. \textit{Hepatol. Commun.} \textbf{10}, e0991 (2026). \url{https://doi.org/10.1097/HC9.0000000000000991}.
\bibitem{himmelstein2017} Himmelstein, D. S. et al. Systematic integration of biomedical knowledge prioritizes drugs for repurposing. \textit{eLife} \textbf{6}, e26726 (2017). \url{https://doi.org/10.7554/eLife.26726}.
\bibitem{ref5} Vittor, L. et al. Unifying biomedical knowledge in a modern multimodal graph. \textit{arXiv} 2604.27269v3 (2026), preprint. \url{https://arxiv.org/abs/2604.27269v3}; \url{https://doi.org/10.48550/arXiv.2604.27269}.
\bibitem{ref4} Hoang, S. A. et al. Gene Expression Predicts Histological Severity and Reveals Distinct Molecular Profiles of Nonalcoholic Fatty Liver Disease. \textit{Sci. Rep.} \textbf{9}, 12541 (2019). \url{https://doi.org/10.1038/s41598-019-48746-5}.
\bibitem{ref7} Ramachandran, P. et al. Resolving the fibrotic niche of human liver cirrhosis at single-cell level. \textit{Nature} \textbf{575}, 512–518 (2019). \url{https://doi.org/10.1038/s41586-019-1631-3}.
\bibitem{ref8} Edgar, R., Domrachev, M. \& Lash, A. E. Gene Expression Omnibus: NCBI gene expression and hybridization array data repository. \textit{Nucleic Acids Res.} \textbf{30}, 207–210 (2002). \url{https://doi.org/10.1093/nar/30.1.207}.
\bibitem{ref10} Fujiwara, N. et al. Molecular signatures of long-term hepatocellular carcinoma risk in nonalcoholic fatty liver disease. \textit{Sci. Transl. Med.} \textbf{14}, eabo4474 (2022). \url{https://doi.org/10.1126/scitranslmed.abo4474}.
\bibitem{ref9} Love, M. I., Huber, W. \& Anders, S. Moderated estimation of fold change and dispersion for RNA-seq data with DESeq2. \textit{Genome Biol.} \textbf{15}, 550 (2014). \url{https://doi.org/10.1186/s13059-014-0550-8}.
\bibitem{ref11} Benjamini, Y. \& Hochberg, Y. Controlling the False Discovery Rate: A Practical and Powerful Approach to Multiple Testing. \textit{J. R. Stat. Soc. B} \textbf{57}, 289–300 (1995). \url{https://doi.org/10.1111/j.2517-6161.1995.tb02031.x}.
\bibitem{ref12} Liberzon, A. et al. The Molecular Signatures Database (MSigDB) hallmark gene set collection. \textit{Cell Syst.} \textbf{1}, 417–425 (2015). \url{https://doi.org/10.1016/j.cels.2015.12.004}.
\bibitem{ref13} Fang, Z., Liu, X. \& Peltz, G. GSEApy: a comprehensive package for performing gene set enrichment analysis in Python. \textit{Bioinformatics} \textbf{39}, btac757 (2023). \url{https://doi.org/10.1093/bioinformatics/btac757}.
\end{thebibliography}



\clearpage
\section*{Supplementary information}
\begingroup
\renewcommand{\thetable}{S\arabic{table}}
\renewcommand{\theHtable}{S\arabic{table}}
\setcounter{table}{0}
\begin{table}[H]\centering
\caption{Cluster-number diagnostics for all evaluated solutions}\label{tab:S1}
\fontsize{9}{11}\selectfont
\begin{tabular}{>{\raggedright\arraybackslash}p{0.045\textwidth}>{\raggedright\arraybackslash}p{0.13\textwidth}>{\raggedright\arraybackslash}p{0.245\textwidth}>{\raggedright\arraybackslash}p{0.245\textwidth}>{\raggedright\arraybackslash}p{0.13\textwidth}}
\toprule
\textbf{k} & \textbf{Silhouette} & \textbf{Gene-subset ARI (mean $\pm$ SD)} & \textbf{Patient-bootstrap ARI (mean $\pm$ SD)} & \textbf{Gene-subset mean ARI $\geq0.95$} \\
\midrule
3 & 0.5874 & 0.9811 $\pm$ 0.0183 & 0.7913 $\pm$ 0.0627 & Yes \\
4 & 0.4378 & 0.9708 $\pm$ 0.0174 & 0.6787 $\pm$ 0.0806 & Yes \\
5 & 0.3522 & 0.9835 $\pm$ 0.0107 & 0.5919 $\pm$ 0.0623 & Yes \\
6 & 0.3167 & 0.8164 $\pm$ 0.1717 & 0.5137 $\pm$ 0.0727 & No \\
7 & 0.2804 & 0.9343 $\pm$ 0.0452 & 0.4495 $\pm$ 0.0478 & No \\
8 & 0.2841 & 0.8676 $\pm$ 0.1104 & 0.4257 $\pm$ 0.0455 & No \\
\bottomrule
\end{tabular}
\par\vspace{6pt}
\begin{minipage}{0.94\textwidth}\small
Gene-subset diagnostics used 50 refits with 80\% of genes; patient-bootstrap diagnostics used a separate 50 refits resampling patients within studies. Values are reproduced from the cluster-number diagnostics, without reclustering. The selection rule chose the largest k with mean gene-subset ARI at least 0.95: k = 5. The highest silhouette score occurred at k = 3; patient-bootstrap ARI was not the selection criterion. Means and SDs are rounded to four decimal places for display; eligibility was assessed using the unrounded means.
\end{minipage}
\end{table}
\endgroup

\subsection*{Supplementary Methods: implementation checks and reference-file inventory}
For transition scoring, the comparison with the largest absolute available full-data coefficient was selected, with ties resolved in the order F0--F1, F1--F2, F2--F3, F3--F4. A missing full-data q-value did not satisfy the substantial-difference condition. Missing matched evidence for the selected comparison, or the absence of a selected comparison for a gene, stopped scoring through an assertion rather than being assigned zero. In the final analysis, all 3,845 fibrosis-associated genes had finite coefficients and q-values for all four comparisons and matched evidence for the selected comparison. The reconstructed P/T/X/F/B scores agreed with the analysis table for all 3,845 genes.

The reference files read by the analysis scripts were \nolinkurl{data/raw/hgnc_complete_set.txt} (file modification date 20 August 2026; most recent gene-record modification date within the file 18 August 2026) and the OptimusKG parquet files in \nolinkurl{data/optimuskg} (file modification date 1 September 2026). SHA-256 checksums for all files are provided in the source manifest, \nolinkurl{source_checksums.txt}. File dates and checksums identify these retained files and do not establish official release dates.

\end{document}
